\documentclass[11pt, a4paper, titlepage, dvipdfmx, uplatex]{jsarticle}

\usepackage{script} % スタイルシート(script.sty)の呼び出し


\dtTitle{要記入：日本語題目}{要記入：英語題目} %\dtTitle{要記入：日本語題目}{要記入：英語題目}
\dtAuthor{要記入：氏名}{要記入：学籍番号}{要記入：入学年} % \dtAuthor{要記入：氏名}{要記入：学籍番号}{要記入：入学年}
\date{要記入：提出日}


\begin{document}
\maketitle

\begin{abstract}
要記入：英文Abstract。本文の目的・方法・結果・結論と一致させる。

\end{abstract}

\newpage

\tableofcontents % 目次の出力
\thispagestyle{empty}
\newpage

\setcounter{page}{1}
\section{はじめに}
要記入：研究資料に基づいて執筆する。


\section{関連研究}
要記入：研究資料に基づいて執筆する。


\section{評価実験}
要記入：研究資料に基づいて執筆する。


\section{結果と考察}
要記入：研究資料に基づいて執筆する。


\section{まとめ}
要記入：研究資料に基づいて執筆する。


\newpage
\section*{謝辞}
%\addcontentsline{toc}{section}{謝辞} %謝辞を目次に表示する
要記入：研究資料に基づいて執筆する。

\newpage

\begin{onehalfspace}
%\addcontentsline{toc}{section}{参考文献} %参考文献を目次に表示する
\bibliographystyle{junsrt}
\bibliography{ref}
\newpage
\end{onehalfspace}

% 本人の研究業績を掲載する場合は、確認済みの文献と別リストを追加する。

\end{document}
